\documentclass[11pt]{article}
\usepackage[utf8]{inputenc}
\usepackage[T1]{fontenc}
\usepackage{amsmath}
\usepackage{amssymb}
\usepackage{multicol}
\usepackage{geometry}
\usepackage{tikz}
\usepackage{enumitem}
\usepackage{xcolor}
\usepackage{titlesec}

% Configurações de layout
\geometry{a4paper, left=1cm, right=1cm, top=0.5cm, bottom=1.2cm}
\setlength{\columnseprule}{0.4pt}

% Cores para títulos
\titleformat{\section}{\normalfont\Large\bfseries\color{blue}}{\thesection}{1em}{}
\titleformat{\subsection}{\normalfont\large\bfseries\color{red}}{\thesubsection}{1em}{}


\title{\textcolor{red}{Revisão: Medidas de Tendência Central \\
    Média, Mediana e Moda}}
\author{Professor: Jefferson}
\date{}

\begin{document}

\maketitle
\vspace{-1.0cm}

\begin{multicols}{2}

\section*{Instruções}
\textbf{Copie e responda no caderno.}

\begin{enumerate}[leftmargin=*]

\item Calcule a média dos conjuntos:
\begin{enumerate}[label=\alph*)]
\item 10, 12, 14, 16, 18
\item 5, 8, 11, 14, 17, 20
\item 2, 4, 6, 8, 10, 12, 14
\item 15, 20, 25, 30, 35
\end{enumerate}

\item Determine a mediana:
\begin{enumerate}[label=\alph*)]
\item 3, 5, 7, 9, 11
\item 2, 4, 6, 8, 10, 12
\item 10, 15, 20, 25, 30, 35, 40
\item 1, 3, 5, 7, 9, 11, 13, 15
\end{enumerate}

\item Identifique a moda:
\begin{enumerate}[label=\alph*)]
\item 2, 3, 4, 4, 5, 6, 6, 6, 7
\item 10, 12, 12, 14, 14, 16, 18
\item 5, 5, 5, 8, 8, 10, 10, 10, 12
\item 1, 2, 3, 4, 5, 6
\end{enumerate}

\item Para o conjunto: 12, 15, 18, 21, 24, 27, 30
\begin{enumerate}[label=\alph*)]
\item Calcule a média
\item Determine a mediana
\item Identifique a moda
\end{enumerate}

\item Um aluno tirou as seguintes notas: 7,0; 8,5; 6,0; 9,0; 7,5. Qual sua média?

\item As idades de um grupo são: 12, 14, 13, 15, 12, 16, 14, 13, 12. Determine:
\begin{enumerate}[label=\alph*)]
\item Média
\item Mediana
\item Moda
\end{enumerate}

\end{enumerate}


\begin{enumerate}[leftmargin=*, start=7]

\item Uma empresa tem 5 funcionários com salários: R\$ 1800, R\$ 2000, R\$ 2200, R\$ 2500, R\$ 5000.
\begin{enumerate}[label=\alph*)]
\item Calcule a média salarial
\item Determine a mediana
\item Qual medida representa melhor o salário típico? Por quê?
\end{enumerate}

\item Em uma prova, as notas foram: 5, 6, 7, 8, 8, 9, 10, 10, 10.
\begin{enumerate}[label=\alph*)]
\item Qual a média?
\item Qual a mediana?
\item Qual a moda?
\end{enumerate}

\item Complete a tabela com as medidas:

\begin{center}
    \small
\begin{tabular}{|c|c|c|c|}
\hline
\textbf{Dados} & \textbf{Média} & \textbf{Mediana} & \textbf{Moda} \\
\hline
4, 8, 12, 16, 20 & & & \\
\hline
3, 5, 5, 7, 9, 11 & & & \\
\hline
10, 20, 30, 40, 50, 60 & & & \\
\hline
\end{tabular}
\end{center}

\item As temperaturas de uma semana: 22, 24, 23, 25, 24, 26, 24 (em °C).
\begin{enumerate}[label=\alph*)]
\item Qual a temperatura média?
\item Qual a temperatura mediana?
\item Qual a moda?
\end{enumerate}

\item Em uma pesquisa sobre número de irmãos:
\[
\begin{array}{c|c}
\text{Nº de irmãos} & \text{Frequência} \\
\hline
0 & 5 \\
1 & 8 \\
2 & 6 \\
3 & 3 \\
4 & 2 \\
\end{array}
\]
Calcule a média, mediana e moda.

\item Uma loja vendeu em 7 dias: 12, 15, 10, 18, 20, 15, 14 produtos.
\begin{enumerate}[label=\alph*)]
\item Quantidade média vendida
\item Quantidade mediana
\item Quantidade modal
\end{enumerate}

\end{enumerate}

\end{multicols}

\end{document}
